\begin{textsample}{POS Dim 3 – persona\_grok – Score 17.00 – t192\_grok.txt}  \label{ex:f3_pos_039}
In the midst of the climate crisis, nuclear power is often touted as a quick, safe, and efficient way to decarbonize our \textbf{energy} \textbf{systems}.
But this promotion overlooks some harsh realities that make it far from the ideal \textbf{solution}.
One major issue is the high \textbf{cost}.
Nuclear \textbf{energy} comes in at \$112 to \$189 per megawatt-hour, which is significantly more expensive than renewables.
On \textbf{top} of that, building nuclear plants takes an \textbf{average} of 10 years, delaying any potential benefits at a \textbf{time} when we \textbf{need} rapid action.
Nuclear facilities are also vulnerable to environmental stresses like heatwaves and water shortages, which can disrupt operations.
And while it 's pitched as low-carbon, the process involves significant CO2 \textbf{emissions} from construction and uranium mining and processing.
Then there 's the \textbf{problem} of radioactive \textbf{waste}, which remains unmanaged and poses long-term risks.
Taxpayers \textbf{end} up footing the bill for \textbf{billions} in \textbf{costs} associated with these projects.
Security concerns add another layer of worry, with risks from terrorism, natural disasters, or conflicts—as we 've seen in Ukraine, where nuclear sites have been caught in the crossfire of war.
Even newer \textbf{technologies}, such as EPR reactors and fusion, are plagued by delays and \textbf{budgets} that spiral out of \textbf{control}.
In contrast, renewables provide faster, cheaper, and more scalable \textbf{options} to achieve the urgent \textbf{emissions} \textbf{reductions} we \textbf{need} right now.

% matched lemmas: average, billion, budget, control, cost, emission, end, energy, need, option, problem, reduction, solution, system, technology, time, top, waste
\end{textsample}
